\section{Related Work} 
\label{sec:related}
In recent literature, training overcomplete respresentations is often advocated as overcomplete bases can transform local minima into saddle points \cite{dauphin2014identifying} or help discover robust solutions \cite{lewicki2000learning}. This indicates that significant model compression is often possible without sacrificing accuracy and that the model's accuracy does not rely
on precise weight values~\cite{keskar2016large}. Most of the recent work on model compression exploits this inherent sparsity of overcomplete represantations. These works include low precision computations~\cite{anwar2015fixed,courbariaux2014low} and quantization of model weights~\cite{han2015deep,zhou2017incremental}. There are also methods which
prune the network by dropping connections with low weights ~\cite{wen2016learning,see2016compression} or use sparse encoding ~\cite{han2015learning}. These methods in practice often suffer from quantization loss especially if the network is deep due to a large number of low-precision multiplications during forward pass. Quantization loss is hard to recover through further training due to the non-trivial nature of backpropagation in low precision~\cite{lin2015neural}. There has been some work on compression aware training~\cite{polino2018model} via model distillation~\cite{hinton2015distilling} that addresses this gap in accuracy. However, both quantized distillation and differentiable quantization methods introduce additional latency due to their slow training process that requires careful tuning, and may not be a good fit for practical systems where a compressed model needs to be rapidly tuned and deployed in production. \\
%We explore a different perspective to approach the compression problem. We argue that since the parameter space of a DNN is often overcomplete, the representation of an input is not an unique combination of basis vectors (Lewicki and Sejnowski 2000). This means that it is possible to find a lower dimensional parameter space which can represent the model with minimal information loss. %Extending this intuition, we use low rank matrix factorization techniques for compressing large weight matrices in deep
%models, focusing on the word embedding matrix which is the size bottleneck for many NLP tasks
 %Rahul:this feels incomplete
Low-rank matrix factorization (LMF) is a very old dimensionality reduction technique widely used in the matrix completion
literature. For example, see \cite{recht2013parallel} and references therein. 
However, there has been  relatively limited work on applying LMF to deep neural models. \cite{sainath2013low,lu2016learning} used low rank matrix factorization for neural speech models.While~\cite{sainath2013low}
reduces parameters of DNN before training,~\cite{xue2013restructuring}
restructures the network using SVD introducing bottleneck layers in the weight
matrices. However, for typical NLP tasks, introducing bottleneck layers between the deep layers of the model does not significantly decrease model size since the large majority of the parameters in NLP models come from the input word embedding matrix.
Building upon this prior work, we apply LMF based ideas to reduce the embedding space of NLP models. Instead of exploiting the sparsity of parameter space we argue that for overcomplete representation the input embedding matrix is not a unique combination of basis vectors~\cite{lewicki2000learning}. Thus, it is possible to find a lower dimensional parameter space which can represent the model inputs with minimal information loss. Based on this intuition, we apply LMF techniques for compressing the word embedding matrix which is the size bottleneck for many NLP tasks. 


%Deep Learning has achieved great success in various NLP tasks over the last few years such as sequence tagging  \cite{chung_rnn}, \cite{ma_hovy_2016} and sentence classification \cite{liu_multi_timescale}, \cite{conv_text_kim} \cite{Socher_recursive_2013}. The NLP literature applies and combines common neural models including LSTMs~\cite{hochreiter_lstm}, CNNs~\cite{lecun98}, DAN~\cite{iyyer2015deep} etc. While traditional machine learning approaches extract hand-designed, task-specific, features which are fed into a shallow model, neural models pass the input through several feedforward, recurrent or convolutional layers of feature extraction that are trained to learn task-specific feature representations.

%Deep neural models often have large memory-footprint which poses significant deployment challenges for practical systems that often have memory and computing power constraints. Recent literature suggests that often DNN based models are over-complete, which makes the model robust by transforming local minima to saddle points ~\cite{dauphin2014identifying}. This indicates that significant model compression is often possible without severe accuracy regression. Related work has explored compression either by compressing an already trained model or by developing compression aware training techniques.

%Works that propose compression after training include low precision computations~\cite{anwar2015fixed,courbariaux2014low} and quantization of model weights~\cite{han2015deep,zhou2017incremental}. There are also methods which prune the network by dropping connections with low weights~\cite{wen2016learning,han2015learning,see2016compression}. A lot of these methods can be used in conjunction with our work as well. Recently there has been work on compression aware training via model distillation~\cite{hinton2015distilling,polino2018model} as well. However, model distillation can be a slow training process that requires careful tuning, and may not be a good fit for practical systems where a compressed model needs to be rapidly tuned and deployed in production.
% [TBA: low precision computations, quantization , distillations , pruning, hashing, bucketing, compact codes, @Anish complete this].  Compression aware training techniques include [TBA: Distillation and Quantization, fake quantization etc, @Anish complete this]

%There has been relatively less work in compressing word embeddings specifically for NLP. ~\cite{shu2017compressing} do hashing on the embeddings while learning the discrete codes by doing the gumbel softmax trick. 
%Low-rank matrix factorization (LMF)\cite{mnih2008probabilistic,sainath2013low,lu2016learning}) has also been widely used in traditional machine learning models as a method for dimensionality reduction. ~\cite{sainath2013low} used low rank matrix factorization for neural speech models. ~\cite{denton2014exploiting}  and more recently ~\cite{kossaifi2017tensor} used similar ideas on contracting tensor layers in CNN. Our work builds upon these spectral projection based compression works on and proposes similar matrix compression techniques for NLP tasks, focusing on compressing the input word embedding matrices and re-training them for downstream tasks.  

 
 
% Related works in NLP also include compressing word embeddings offline before tuning them on downstream tasks [** Compressing Embedding - Existing work @Anish complete this].

%TODO: Add this somewhere in the last paragraph, when discussing NLP work.
%However, even for small NLP tasks this embedding layer can be large.  For example, to represent a vocabulary of 100K words using 300 dimensional Glove embeddings the embedding matrix would have to hold 60M parameters. Even for a simple sentiment analysis model the embedding parameters are 98.8\% of the total network ~\cite{shu2017compressing}. 




%[@Anish here you can optionally discuss some disadvantages of existing methods that our method addresses]

%As discussed in \ref{sec:related} there are two related but slightly differnet research directions. 
%One set of work tries to compress already-trained models eg. quantization. These approaches often introduce compression loss resulting in lower accuracy. However, these often have faster inference time. In an attempt to regain this accuracy loss quantization 
%The other direction is to train compressed networks. However, it is not always trivial to train quantized or hashed networks since 
%Existing works on compressing DNNs usually introduce some loss in accuracy. 
%Some recent works that are able to achieve almost lossless compression are usually very complex and increase training time considerably. 


%Recognizing the model size problems for deep learning various approaches have been proposed in the past, including the hashing trick \cite{chen2015compressing}, quantization \cite{hubara2016quantized} and combination of pruning, trained quantization and Huffman encoding \cite{han2015deep}.
 









